\section{Các phương pháp xây dựng mô hình bảo toàn tính đối xứng}
\label{sec:phuong-phap}

% ============================================================
% MỤC TIÊU CHƯƠNG:
% Trình bày 2 hướng chính để đưa tính đối xứng vào AI,
% cùng các ứng dụng thực tế và giới hạn của chúng.
% Nguồn: Slides Part 1 (trang 15–55), Tutorial Part 1
% ============================================================

% ─────────────────────────────────────────────────────────────
% HƯỚNG 1: CAN THIỆP DỮ LIỆU (OFF-THE-SHELF MODELS)
% Dùng model bình thường + xử lý dữ liệu đầu vào/đầu ra
% để đạt được tính invariance/equivariance.
% ─────────────────────────────────────────────────────────────

\subsection{Hướng 1: Can thiệp dữ liệu trên các mô hình có sẵn}
\label{subsec:off-the-shelf}

Trước khi đi vào các phương pháp xây dựng kiến trúc hoàn toàn mới, một hướng tiếp cận tự nhiên hơn là tận dụng lại những mô hình học máy đã được huấn luyện tốt (thứ mà cộng đồng nghiên cứu đã đổ rất nhiều công sức phát triển) rồi can thiệp vào phía dữ liệu để ép mô hình phải tôn trọng tính đối xứng. Hướng này được gọi là hướng can thiệp dữ liệu, hay còn được gọi là cách tiếp cận với mô hình hộp đen.

Ưu điểm cốt lõi của hướng này là mô hình cơ sở (dù là Transformer, CNN hay bất kỳ kiến trúc nào khác) được giữ nguyên hoàn toàn, kéo theo toàn bộ sức mạnh biểu diễn, kỹ thuật tối ưu, và quy trình huấn luyện của nó vẫn còn hiệu lực. Điều đó có nghĩa là người dùng không cần thiết kế lại bất cứ điều gì từ đầu. Ngoài ra, nếu mô hình cơ sở đã có sẵn một số tính chất đối xứng nhất định, việc kết hợp thêm các nhóm đối xứng khác thông qua can thiệp dữ liệu là hoàn toàn khả thi mà không gây mâu thuẫn.

Trong hướng này có ba kỹ thuật chính: tăng cường dữ liệu, chuẩn hóa, và lấy trung bình nhóm/khung. Ba kỹ thuật này nằm trên một phổ liên tục: từ một phép biến đổi duy nhất ở một đầu, đến lấy trung bình trên toàn bộ nhóm $G$ ở đầu kia, với lấy trung bình khung nằm ở giữa như một điểm cân bằng giữa chi phí tính toán và chất lượng đảm bảo tính đối xứng.

\subsubsection{Tăng cường dữ liệu}
\label{subsubsec:data-augmentation}

\paragraph{Ý tưởng cơ bản.}
Tăng cường dữ liệu là kỹ thuật đơn giản và phổ biến nhất để đưa tính đối xứng vào mô hình: với mỗi điểm dữ liệu $x$ trong tập huấn luyện, ta sinh thêm các biến thể $g \cdot x$ cho nhiều phần tử $g \in G$ khác nhau và đưa chúng vào tập huấn luyện cùng nhãn gốc. Khi tập huấn luyện bao phủ đủ nhiều biến thể của quỹ đạo $\mathrm{Orb}_G(x)$, mô hình sẽ dần học được rằng phép biến đổi $g$ không làm thay đổi nhãn, tức là dần hình thành tính bất biến.

Kỹ thuật này cực kỳ linh hoạt vì nó có thể áp dụng cho cả học có giám sát lẫn học tự giám sát. Trong học tương phản, người ta thường tạo cặp dữ liệu $(x, g \cdot x)$ và đặt chúng ở vị trí ``dương'' (positive pair), buộc mô hình phải đặt cả hai vào cùng một vùng trong không gian ẩn. Qua đó, toàn bộ quỹ đạo của $x$ dưới tác động của $G$ sẽ được ánh xạ về một biểu diễn duy nhất.

\paragraph{Giới hạn: chỉ đảm bảo tính đối xứng xấp xỉ.}
Dù đơn giản và hiệu quả trong thực tế, tăng cường dữ liệu mang một hạn chế quan trọng về mặt lý thuyết: nó không đảm bảo tính bất biến ở mọi điểm. Sau khi huấn luyện, không có gì chắc chắn rằng với một điểm dữ liệu mới $x'$ chưa thấy, mô hình sẽ cho ra cùng một kết quả cho $x'$ và $g \cdot x'$. Mô hình có thể đã \textit{học ước lượng} tính bất biến trong vùng lân cận của dữ liệu huấn luyện, nhưng không có gì đảm bảo điều đó khái quát hóa ra ngoài phân phối.

\begin{warningbox}{Tăng cường dữ liệu chỉ đưa ra xấp xỉ, không phải bảo đảm}
Biểu hiện thực tế của hạn chế này: mô hình huấn luyện với tăng cường dữ liệu thường hoạt động tốt trên bộ test trong phân phối, nhưng rất dễ bị đánh lừa khi gặp các đầu vào bị biến đổi ở góc độ chưa từng xuất hiện trong tập huấn luyện. Đây không phải lỗi do lượng dữ liệu ít, mà là một giới hạn căn bản: không có đảm bảo tính bất biến ở mọi nơi.
\end{warningbox}

\paragraph{Mở rộng: học biểu diễn equivariant.}
Một bước phát triển thú vị từ tăng cường dữ liệu là học biểu diễn \textit{equivariant} thay vì chỉ invariant. Trong bài toán học biểu diễn, đôi khi ta không muốn hoàn toàn ``quên đi'' phép biến đổi mà dữ liệu đã trải qua; ta muốn không gian ẩn ghi nhớ lại sự biến đổi đó theo một cách có cấu trúc. Cụ thể, nếu $a$ là một phép tăng cường và $f$ là mạng mã hóa, thì mục tiêu là:
\begin{equation}
    f(a(x)) = T_a \cdot f(x),
\end{equation}
trong đó $T_a$ là một phép biến đổi tác động lên không gian ẩn, tương ứng với $a$. Đây chính là định nghĩa equivariance áp dụng vào bài toán học biểu diễn. Hình~\ref{fig:equivariant-repr} minh họa trực quan ý tưởng này.

\begin{figure}[H]
    \centering
    \begin{tikzpicture}[
        node distance=2.2cm and 3.8cm
        every node/.style={font=\normalsize},
        box/.style={draw, rounded corners=4pt, minimum width=1.5cm, minimum height=0.8cm,,
                    fill=blue!8, thick},
        arr/.style={->, >=Stealth, thick, blue!60!black}
    ]
        \node[box] (x)  {$x$};
        \node[box] (ax) [right=of x]  {$a(x)$};
        \node[box] (fx) [below=of x]  {$f(x)$};
        \node[box] (fax)[below=of ax] {$f(a(x))$};

        \draw[arr] (x)  -- node[above]{augmentation $a$} (ax);
        \draw[arr] (x)  -- node[left] {encoder $f$}      (fx);
        \draw[arr] (ax) -- node[right]{encoder $f$}      (fax);
        \draw[arr] (fx) -- node[below]{transform $T_a$}  (fax);

        \node[font=\footnotesize, align=center] at (1.9,-1.1)
            {\textit{both paths}\\[2pt]\textit{yield the same result}};
    \end{tikzpicture}
    \caption{Sơ đồ giao hoán của học biểu diễn equivariant: áp dụng phép tăng cường $a$ rồi mã hóa, hoặc mã hóa rồi tác động $T_a$ trong không gian ẩn, đều cho cùng kết quả.}
    \label{fig:equivariant-repr}
\end{figure}

Một cách tiếp cận trực tiếp là dùng một mạng nhỏ $\mathrm{MLP}(a, f(x))$ để dự đoán $f(a(x))$, từ đó đặt hàm mất mát:
\begin{equation}
    \mathcal{L} = \bigl\| f(a(x)) - \mathrm{MLP}(a,\, f(x)) \bigr\|.
\end{equation}
Tuy nhiên, cách này có một nhược điểm: không có gì đảm bảo $T_a$ sẽ tương thích với phép kết hợp nhóm, tức là $T_{a_2 \circ a_1} \neq T_{a_2} \circ T_{a_1}$ có thể xảy ra.

\paragraph{Phiên bản tuyến tính: ràng buộc bảo toàn tích vô hướng.}
Để khắc phục hạn chế đó, các nghiên cứu sau này đề xuất một cách đơn giản hơn: giả định rằng $T_a$ là một phép quay tác động lên không gian ẩn, tức là $T_a = R_a$ là ma trận trực giao. Khi đó, ta không cần trực tiếp học $R_a$, mà thay vào đó tận dụng tính chất của phép quay: nó bảo toàn tích vô hướng giữa hai vector. Cụ thể, nếu $f(a(x'))^\top f(a(x)) = f(x')^\top f(x)$ với mọi cặp $(x, x')$ và phép tăng cường $a$, thì $T_a$ nhất thiết phải là một phép quay. Hàm mất mát tương ứng là:
\begin{equation}
    \mathbb{E}_{a,\, x,\, x'} \left[ f(a(x'))^\top f(a(x)) - f(x')^\top f(x) \right]^2 = 0.
\end{equation}

Khi hàm mất mát bằng 0, mỗi $T_a = R_a$ là một phép quay trong không gian ẩn và tự động tương thích với phép kết hợp nhóm --- tức là $\{T_a\}_{a \in G}$ tạo thành một biểu diễn nhóm của $G$. Điều đáng chú ý là phép tăng cường $a$ tác động lên đầu vào không nhất thiết phải là một phép quay; đây là hệ quả của định lý Peter--Weyl trong lý thuyết nhóm. Hướng này được phát triển bởi Shakerinava et al.\ (2022) và Gupta et al.\ (2023).

\subsubsection{Chuẩn hóa}
\label{subsubsec:canonicalization}

\paragraph{Ý tưởng cơ bản.}
Kỹ thuật chuẩn hóa theo đuổi một chiến lược hoàn toàn khác: thay vì làm phong phú tập dữ liệu bằng nhiều biến thể, ta chuyển mỗi điểm dữ liệu $x$ về một dạng chuẩn duy nhất trước khi đưa vào mô hình. Miễn là hàm chuẩn hóa $c(\cdot)$ ánh xạ mọi phần tử trong cùng một quỹ đạo về cùng một điểm duy nhất, thì mô hình sẽ tự động bất biến:
\begin{equation}
    c(g \cdot x) = c(x) \quad \forall\, g \in G.
\end{equation}

Ví dụ điển hình nhất là sắp xếp cho bài toán invariant với hoán vị: với bất kỳ chuỗi nào, ta sắp xếp nó theo thứ tự tăng dần trước khi đưa vào mô hình. Vì mọi hoán vị của cùng một chuỗi đều cho ra cùng một dạng đã sắp xếp, mô hình luôn nhận cùng một đầu vào và cho ra cùng một kết quả. Để đạt equivariance thay vì chỉ invariance, ta thực hiện thêm một bước: sau khi mô hình xử lý dạng chuẩn, ta áp dụng phép biến đổi ngược để khôi phục thứ tự ban đầu. Hình~\ref{fig:canon-pipeline} minh họa toàn bộ quy trình này.

\begin{figure}[H]
    \centering
    \begin{tikzpicture}[
        node distance=0.5cm and 1.4cm,
        every node/.style={font=\small},
        box/.style={draw, rounded corners=4pt, minimum width=2.0cm, minimum height=0.75cm,
                    fill=orange!10, thick, align=center},
        op/.style={draw, circle, minimum size=0.7cm, fill=green!15, thick, font=\footnotesize},
        arr/.style={->, >=Stealth, thick}
    ]
        \node[box] (x)      {$x$\\(đầu vào)};
        \node[op]  (c)      [right=of x]          {$c(\cdot)$};
        \node[box] (cx)     [right=of c]          {$c(x)$\\(dạng chuẩn)};
        \node[op]  (f)      [right=of cx]         {$f$};
        \node[box] (fcx)    [right=of f]          {$f(c(x))$\\(đầu ra $f$)};
        \node[op]  (inv)    [right=of fcx]        {$c^{-1}$};
        \node[box] (out)    [right=of inv]        {đầu ra\\equivariant};

        \draw[arr] (x)   -- (c);
        \draw[arr] (c)   -- (cx);
        \draw[arr] (cx)  -- (f);
        \draw[arr] (f)   -- (fcx);
        \draw[arr] (fcx) -- (inv);
        \draw[arr] (inv) -- (out);

        % Invariant path (shorter)
        \draw[arr, dashed, gray] (fcx.south) -- ++(0,-0.6)
            node[below, align=center, font=\footnotesize, gray]
            {nếu chỉ cần invariance:\\dừng ở đây};
    \end{tikzpicture}
    \caption{Pipeline chuẩn hóa: đầu vào $x$ được ánh xạ về dạng chuẩn $c(x)$, xử lý bằng mô hình $f$, sau đó áp dụng phép biến đổi ngược $c^{-1}$ để đạt equivariance. Nếu chỉ cần invariance thì bỏ qua bước cuối.}
    \label{fig:canon-pipeline}
\end{figure}

Trong ví dụ sắp xếp, nếu dãy $[5,\, 3,\, 4]$ được sắp thành $[3,\, 4,\, 5]$, mô hình xử lý và cho ra kết quả theo thứ tự đó, rồi ta hoán vị ngược lại để đầu ra có thứ tự tương ứng với đầu vào gốc.

Hàm chuẩn hóa $c(\cdot)$ không nhất thiết phải được thiết kế bằng tay; nó cũng có thể được học từ dữ liệu thông qua một mạng neural phụ trợ \cite{kaba2022equivariance}. Hướng này mở ra khả năng tận dụng sức mạnh của deep learning để tự động tìm ra cách biểu diễn chuẩn phù hợp nhất cho bài toán cụ thể.

\paragraph{Thách thức 1: không phải lúc nào cũng dễ xác định dạng chuẩn.}
Đối với chuỗi số, sắp xếp là một lựa chọn tự nhiên và đơn giản. Nhưng không phải mọi loại dữ liệu đều có một dạng chuẩn rõ ràng. Hãy lấy ví dụ với đồ thị: một đồ thị $n$ nút có thể được mô tả bởi $n!$ ma trận kề khác nhau, tương ứng với $n!$ cách đánh số các nút. Tìm một biểu diễn chuẩn chung cho tất cả các cách đánh số đó về cơ bản tương đương với việc giải bài toán kiểm tra đẳng cấu đồ thị (graph isomorphism) --- một bài toán đã được nghiên cứu trong nhiều thập kỷ nhưng vẫn chưa có thuật toán đa thức được chứng minh. Tương tự với đám mây điểm trong không gian 3D kết hợp đối xứng quay và hoán vị, việc chuẩn hóa trở nên rất phức tạp.

\paragraph{Thách thức 2: vấn đề tính liên tục.}
Đây là hạn chế lý thuyết nghiêm trọng nhất của kỹ thuật chuẩn hóa. Tính liên tục yêu cầu rằng nếu hai đầu vào gần nhau thì hai đầu ra cũng phải gần nhau. Với chuỗi số vô hướng, sắp xếp duy trì tính liên tục vì một thay đổi nhỏ thường không làm đảo lộn toàn bộ thứ tự. Tuy nhiên, câu chuyện hoàn toàn khác khi dữ liệu là các vector đa chiều.

Hãy xem xét ví dụ cụ thể từ slides của khóa học (trang 14 trong file slides NeurIPS 2025): ta có hai đám mây điểm gần nhau, mỗi đám gồm hai vector. Đám thứ nhất chứa $[5,\, 3,\, 3]^\top$ và $[1,\, 7,\, 8]^\top$; đám thứ hai chứa $[5,\, 3.01,\, 3]^\top$ và $[1,\, 7,\, 8]^\top$ --- chỉ thay đổi một phần tử nhỏ từ $3$ thành $3.01$. Nếu ta sắp xếp các vector theo thứ tự từ điển, thì thứ tự sắp xếp của hai đám sẽ khác nhau hoàn toàn vì $3.01$ khiến thứ tự tương đối của hai chiều thứ hai và thứ ba bị đảo ngược. Kết quả là hai đầu vào cực kỳ gần nhau lại được ánh xạ đến hai dạng chuẩn rất xa nhau về khoảng cách Euclidean.
\begin{figure}[H]
    \centering
    \includegraphics[width=0.85\textwidth, keepaspectratio]{graphics/canon_continuity.png}
    \caption{Minh họa vấn đề tính liên tục trong chuẩn hóa: hai đám mây điểm chỉ khác nhau rất nhỏ (số 3 so với 3.01) nhưng sau khi sắp xếp lại cho ra hai dạng chuẩn xa nhau về khoảng cách Euclidean. Nguồn: \cite{jegelka2025tutorial}, trang 14.}
    \label{fig:canon-continuity}
\end{figure}

Kết quả lý thuyết của Dym et al.\ (2024) và Zhang et al.\ (2020) chỉ ra rằng đây không phải vấn đề của riêng một cách sắp xếp cụ thể nào, mà là một giới hạn căn bản không thể tránh được:

\begin{theorem}[Không tồn tại chuẩn hóa liên tục, Dym et al.\ 2024]
\label{thm:no-continuous-canon}
Không tồn tại hàm chuẩn hóa liên tục nào cho nhóm hoán vị $S_n$ tác động lên $\mathbb{R}^{d \times n}$ với $n, d > 1$. Tương tự, không tồn tại frame hữu hạn liên tục nào cho $SO(2)$ trong trường hợp tổng quát.
\end{theorem}

Hệ quả thực tiễn của định lý này là: dù ta thiết kế hàm chuẩn hóa khéo léo đến đâu, sẽ luôn tồn tại những cặp đầu vào gần nhau nhưng có dạng chuẩn cách xa nhau, dẫn đến hành vi không ổn định của mô hình. Đây là lý do frame averaging --- sẽ được trình bày ở mục tiếp theo --- trở thành một lựa chọn hấp dẫn hơn: bằng cách lấy trung bình trên nhiều phép biến đổi thay vì chỉ chọn một, frame averaging khắc phục được vấn đề liên tục này.

\paragraph{Tổng kết về chuẩn hóa.}
Nhìn chung, chuẩn hóa là kỹ thuật hiệu quả và rẻ về tính toán so với lấy trung bình nhóm, đặc biệt khi nhóm $G$ lớn. Nó phù hợp nhất trong các trường hợp dạng chuẩn tự nhiên và dễ xác định (như sắp xếp cho hoán vị trên vô hướng), và cần thận trọng với các bài toán yêu cầu tính ổn định cao trước các nhiễu nhỏ. Bảng~\ref{tab:symmetrization-compare} ở cuối mục này tóm tắt so sánh ba kỹ thuật.

\subsubsection{Lấy trung bình Nhóm và Khung (Group/Frame Averaging)}
\label{subsubsec:frame-averaging}
% ĐÂY LÀ PHẦN CỦA SỸ — Paper 3: Frame Averaging (Puny et al 2022)

% Group Averaging (Reynolds operator):
%   - Invariant: f̄(x) = (1/|G|) Σ_{g∈G} f(g^{-1}·x)
%   - Equivariant: f̄(x) = (1/|G|) Σ_{g∈G} g·f(g^{-1}·x)
%   - Universal approximation (Yarotsky 2021)
%   - Ví dụ: Janossy pooling (Murphy et al 2019),
%     SignNet cho sign flips của eigenvectors (Lim et al 2023)
%   - Nhược điểm: expensive khi |G| lớn → cần approximation

% Frame Averaging (Puny et al 2022):
%   - Thay vì average trên toàn bộ G, chỉ average trên F(X) ⊂ G
%   - F(X) được trích xuất từ data, phải thỏa F(g·X) = g·F(X)
%   - Công thức: ⟨Φ⟩_F(X) = (1/|F(X)|) Σ_{g∈F(X)} ρ₂(g)Φ(ρ₁(g)^{-1}X)
%   - Ví dụ PCA Frame cho point clouds/proteins:
%     F(X) = {(α₁v₁,...,αdvd, t) | αᵢ ∈ {-1,+1}} (sign flips của eigenvectors)
%     Đạt E(3)/SE(3) equivariance với chi phí thấp

% Bảng so sánh 3 phương pháp:
%   Canonicalization:  cheap, NOT continuous
%   Frame Averaging:   cheaper than full, continuous
%   Group Averaging:   most expensive, continuous, universal

% Nội dung đã có trong file gốc (công thức LaTeX của Văn Sỹ)
% — giữ nguyên và mở rộng phần này.

Một trong những phương pháp mạnh mẽ nhất để trang bị tính bất biến (invariance) hoặc tương biến (equivariance) cho một mô hình học máy thông thường là sử dụng toán tử Reynolds, hay còn gọi là kỹ thuật lấy trung bình nhóm (Group Averaging). Cụ thể, cho một mạng nơ-ron cơ sở $\Phi$, ta có thể tạo ra một hàm tương biến $\langle \Phi \rangle_G$ bằng cách tính trung bình đầu ra của $\Phi$ trên toàn bộ các phép biến đổi $g$ thuộc nhóm đối xứng $G$:
\begin{equation}
    \langle \Phi \rangle_G(X) = \frac{1}{|G|} \sum_{g \in G} \rho_2(g)\Phi(\rho_1(g)^{-1}X)
\end{equation}
Mặc dù phương pháp này đảm bảo tính đối xứng tuyệt đối và bảo toàn khả năng biểu diễn của mô hình cơ sở, nó vấp phải trở ngại cực lớn về độ phức tạp tính toán khi nhóm $G$ quá lớn hoặc có số phần tử vô hạn (ví dụ: nhóm hoán vị $S_n$ hoặc nhóm phép quay $SO(3)$).

Để giải quyết bài toán tính toán này, kỹ thuật \textbf{Frame Averaging (Trung bình hóa Khung)} đã được đề xuất \cite{puny2022frame}. Thay vì tính trung bình trên toàn bộ nhóm $G$, phương pháp này chỉ thực hiện lấy trung bình trên một tập hợp con nhỏ $\mathcal{F}(X) \subset G$ được gọi là ``khung'' (frame). Tập hợp $\mathcal{F}(X)$ này được trích xuất dựa trên chính đặc trưng của dữ liệu đầu vào $X$. Để đảm bảo tính tương biến, khung $\mathcal{F}$ phải thỏa mãn điều kiện $\mathcal{F}(\rho_1(g)X) = g\mathcal{F}(X)$. Khi đó, công thức xấp xỉ trở thành:
\begin{equation}
    \langle \Phi \rangle_{\mathcal{F}}(X) = \frac{1}{|\mathcal{F}(X)|} \sum_{g \in \mathcal{F}(X)} \rho_2(g)\Phi(\rho_1(g)^{-1}X)
\end{equation}

\textbf{Ứng dụng PCA Frame cho Đám mây điểm (Point Clouds):}
Trong thực tiễn, đối với dữ liệu không gian 3 chiều như đồ thị phân tử hoặc đám mây điểm, Phân tích thành phần chính (PCA) thường được sử dụng để xây dựng khung. Thuật toán tính tâm của đám mây điểm và trích xuất các vector riêng (eigenvectors) từ ma trận hiệp phương sai. Khi dữ liệu đầu vào bị xoay hoặc tịnh tiến, PCA Frame tự động tính phép xoay ngược để chiếu dữ liệu về canonical frame trước khi đưa qua mạng. Phương pháp này đạt tính bất biến với $E(3)$ với chi phí tính toán cực kỳ thấp.

\subsection{Hướng 2: Xây dựng kiến trúc tương biến từ đầu}
\label{subsec:constructed}

Nếu Hướng 1 xem mô hình học máy như một hộp đen và truyền đạt tính đối xứng từ bên ngoài thông qua thao tác dữ liệu, thì Hướng 2 lựa chọn con đường triệt để hơn: mã hóa cứng tính đối xứng vào chính kiến trúc của mạng. Mô hình tương biến được xây dựng theo cấu trúc tổng quát sau.

\begin{figure}[H]
\centering
\begin{tikzpicture}[>=Stealth, thick]
    % Khai báo các node với khoảng cách mép cố định (sử dụng thư viện positioning)
    \node[draw, rounded corners, fill=blue!10, minimum width=1.2cm, minimum height=0.8cm] (x) {$x$};
    \node[draw, rounded corners, fill=orange!20, minimum height=0.8cm, right=0.6cm of x] (L1) {Lớp tương biến $L_1$};
    \node[draw, rounded corners, fill=gray!15, minimum width=1.5cm, minimum height=0.8cm, right=0.6cm of L1] (sig) {$\sigma(\cdot)$};
    \node[draw, rounded corners, fill=orange!20, minimum height=0.8cm, right=0.6cm of sig] (L2) {Lớp tương biến $L_2$};
    \node[draw, rounded corners, fill=green!20, minimum height=0.8cm, right=0.6cm of L2] (pool) {Gộp bất biến};
    \node[draw, rounded corners, fill=red!10, minimum width=1.2cm, minimum height=0.8cm, right=0.6cm of pool] (y) {$y$};

    % Vẽ mũi tên
    \draw[->] (x) -- (L1);
    \draw[->] (L1) -- (sig);
    \draw[->] (sig) -- (L2);
    \draw[->] (L2) -- (pool);
    \draw[->] (pool) -- (y);

    % Ghi chú bên dưới các khối
    \node[below=0.2cm of L1, font=\small, text=orange!70!black] {tương biến};
    \node[below=0.2cm of L2, font=\small, text=orange!70!black] {tương biến};
    \node[below=0.2cm of pool, font=\small, text=green!60!black] {bất biến};
\end{tikzpicture}
\caption{Cấu trúc tổng quát của mô hình tương biến: xen kẽ các lớp tuyến tính tương biến với hàm kích hoạt phi tuyến $\sigma$, kết thúc bằng một lớp gộp bất biến toàn cục.}
\label{fig:equivariant-pipeline}
\end{figure}

\textbf{Ưu điểm cốt lõi của hướng này:}
\begin{itemize}
    \item Tính đối xứng được đảm bảo tuyệt đối tại mọi điểm dữ liệu, kể cả những điểm chưa thấy trong quá trình huấn luyện.
    \item Số tham số giảm mạnh do cơ chế chia sẻ, dẫn đến hiệu quả dữ liệu cao hơn.
    \item Đặc biệt phù hợp khi tính đối xứng là điều chắc chắn về mặt vật lý hoặc toán học (phân tử, robot, đồ thị).
\end{itemize}

% ──────────────────────────────────────────────────────────────

\subsubsection{Lớp tuyến tính tương biến và cơ chế chia sẻ tham số}
\label{subsubsec:linear-equivariant}

\textbf{Nguyên lý xây dựng.} Một lớp tuyến tính $\mathbf{W} \in \mathbb{R}^{n \times n}$ đảm bảo tính tương biến với nhóm $G$ khi và chỉ khi nó thỏa mãn ràng buộc sau với mọi phần tử $\tau \in G$:

\begin{equation}
    P_\tau \mathbf{W} P_\tau^\top = \mathbf{W}
    \label{eq:equivariance-constraint}
\end{equation}

Ràng buộc này kéo theo một hệ quả quan trọng: nếu tồn tại phép biến đổi $\tau \in G$ ánh xạ cặp chỉ số $(i,j)$ thành $(k,l)$, thì bắt buộc $W_{ij} = W_{kl}$. Đây là cơ chế \textbf{chia sẻ tham số}: các vị trí cùng quỹ đạo trong $G$ phải mang cùng giá trị, và mỗi quỹ đạo chỉ tương ứng với một tham số tự do duy nhất.

\vspace{0.5em}
\textbf{Ví dụ 1: Nhóm dịch chuyển tuần hoàn dẫn đến phép tích chập.}

Xét tín hiệu một chiều với nhóm dịch chuyển $G = C_n = \{t_0, t_1, \ldots, t_{n-1}\}$, trong đó $t_i(j) = i + j \pmod{n}$.

\begin{figure}[H]
\centering
\begin{tikzpicture}[scale=0.75]
    % Ma trận W dạng circulant
    \def\n{5}
    \foreach \i in {0,...,4} {
        \foreach \j in {0,...,4} {
            \pgfmathtruncatemacro{\val}{mod(\i-\j+5,5)}
            \pgfmathsetmacro{\clr}{
                ifthenelse(\val==0,"blue!30",
                ifthenelse(\val==1,"red!30",
                ifthenelse(\val==2,"green!30",
                ifthenelse(\val==3,"orange!30","purple!30"))))
            }
            \fill[\clr] (\j*0.7, -\i*0.7) rectangle (\j*0.7+0.65, -\i*0.7+0.65);
            \draw (\j*0.7, -\i*0.7) rectangle (\j*0.7+0.65, -\i*0.7+0.65);
        }
    }
    % Nhãn các màu
    \node[right=0.5cm] at (5*0.7, -0.5*0.7) {\small $w_0$ \textcolor{blue!70!black}{$\blacksquare$}};
    \node[right=0.5cm] at (5*0.7, -1.5*0.7) {\small $w_1$ \textcolor{red!70!black}{$\blacksquare$}};
    \node[right=0.5cm] at (5*0.7, -2.5*0.7) {\small $w_2$ \textcolor{green!70!black}{$\blacksquare$}};
    \node[right=0.5cm] at (5*0.7, -3.5*0.7) {\small $w_3$ \textcolor{orange!70!black}{$\blacksquare$}};
    \node[right=0.5cm] at (5*0.7, -4.5*0.7) {\small $w_4$ \textcolor{purple!70!black}{$\blacksquare$}};
    \node[below=0.3cm] at (2*0.7, -5*0.7) {\small Ma trận $\mathbf{W}$ sau khi áp ràng buộc: chỉ $n=5$ tham số};
\end{tikzpicture}
\caption{Ma trận trọng số $\mathbf{W}$ sau khi áp ràng buộc tương biến với nhóm dịch chuyển tuần hoàn $C_5$. Các ô cùng màu chia sẻ cùng một giá trị tham số. Cấu trúc này chính xác là ma trận circulant của phép tích chập.}
\label{fig:circulant-matrix}
\end{figure}

\begin{itemize}
    \item Số tham số giảm từ $n^2$ xuống còn $n$, tiết kiệm $n$ lần.
    \item Cấu trúc lặp lại theo đường chéo này chính là định nghĩa của phép \textbf{tích chập} (convolution).
    \item \textit{Kết luận:} mạng tích chập không phải là thiết kế tùy hứng mà là hệ quả tất yếu của việc áp đặt tính tương biến với nhóm dịch chuyển.
\end{itemize}

\vspace{0.5em}
\textbf{Ví dụ 2: Nhóm hoán vị toàn phần dẫn đến kiến trúc DeepSets.}

Xét đầu vào là tập hợp $n$ vector $\{x_1, \ldots, x_n\} \subset \mathbb{R}^d$ và ta muốn mô hình bất biến với mọi phép hoán vị phần tử, tức $G = S_n$.

\begin{figure}[H]
\centering
\begin{tikzpicture}[scale=0.75]
    \def\n{5}
    \foreach \i in {0,...,4} {
        \foreach \j in {0,...,4} {
            \pgfmathtruncatemacro{\same}{ifthenelse(\i==\j,1,0)}
            \ifnum\same=1
                \fill[blue!40] (\j*0.7, -\i*0.7) rectangle (\j*0.7+0.65, -\i*0.7+0.65);
            \else
                \fill[orange!30] (\j*0.7, -\i*0.7) rectangle (\j*0.7+0.65, -\i*0.7+0.65);
            \fi
            \draw (\j*0.7, -\i*0.7) rectangle (\j*0.7+0.65, -\i*0.7+0.65);
        }
    }
    \node[right=0.4cm] at (5*0.7, -0.5*0.7) {\small $\alpha_1$ \textcolor{blue!70!black}{$\blacksquare$} (đường chéo)};
    \node[right=0.4cm] at (5*0.7, -1.5*0.7) {\small $\alpha_2$ \textcolor{orange!70!black}{$\blacksquare$} (ngoài đường chéo)};
    \node[below=0.3cm] at (2*0.7, -5*0.7) {\small Ma trận $\mathbf{W}$ với $G = S_5$: chỉ \textbf{2 tham số}};
\end{tikzpicture}
\caption{Ma trận trọng số $\mathbf{W}$ sau khi áp ràng buộc tương biến với nhóm hoán vị $S_n$. Toàn bộ ma trận $n \times n$ chỉ còn 2 giá trị: $\alpha_1$ trên đường chéo và $\alpha_2$ ở tất cả vị trí còn lại.}
\label{fig:permutation-matrix}
\end{figure}

Kết quả là lớp tuyến tính tương biến có dạng cực kỳ đơn giản:

\begin{equation}
    [\mathbf{F}(\mathbf{X})]_i = \alpha_1 x_i + \alpha_2 \sum_{j \neq i} x_j
    \label{eq:deepset-layer}
\end{equation}

\begin{itemize}
    \item Số tham số giảm từ $n^2$ xuống còn $\mathbf{2}$, bất kể kích thước $n$.
    \item Mỗi điểm được biến đổi kết hợp giữa chính nó và tổng tất cả điểm còn lại.
    \item Đây chính xác là kiến trúc \textbf{DeepSets} và \textbf{PointNet}.
    \item Mở rộng: với tensor bậc cao hơn hoặc các nhóm phức tạp hơn (ví dụ hoán vị ma trận kề của đồ thị), số phần tử cơ sở tăng lên nhưng vẫn ít hơn rất nhiều so với $n^2$. Với hoán vị đồng thời hàng và cột của ma trận kề, số cơ sở chỉ là 15.
\end{itemize}

% ──────────────────────────────────────────────────────────────

\subsubsection{Tích chập nhóm}
\label{subsubsec:group-convolution}

\textbf{Từ tích chập thông thường đến tích chập nhóm.}

Tích chập thông thường trên lưới $\mathbb{Z}^2$ thực hiện phép trượt bộ lọc $\psi$ qua toàn bộ các phép tịnh tiến:
\begin{equation}
    (f \star \psi)(x) := \sum_{y \in \mathbb{Z}^2} f(y)\, \psi(y - x)
    \label{eq:regular-conv}
\end{equation}

Quan sát rằng phép trừ $y - x$ chính là tác động của phép tịnh tiến ngược. Điều này gợi ý cách tổng quát hóa tự nhiên: thay toàn bộ phép tịnh tiến bằng một nhóm đối xứng bất kỳ $G$.

\begin{figure}[H]
\centering
\begin{tikzpicture}[>=Stealth, thick]
    % Khối 1: CNN thông thường (Gộp toàn bộ text vào 1 node)
    \node[draw, rounded corners, fill=blue!10, align=center, minimum width=3.5cm, minimum height=1.5cm, inner sep=0.2cm] (cnn) {
        Bộ lọc trượt theo\\
        phép \textbf{tịnh tiến}\\[1ex]
        \textcolor{blue!70!black}{\footnotesize Tương biến tịnh tiến}
    };
    \node[above=0.2cm of cnn, font=\small\bfseries] {CNN thông thường};

    % Khối Mũi tên ở giữa
    \node[right=1.2cm of cnn, align=center] (arrow) {\Large $\Rightarrow$\\[0.5ex] \footnotesize tổng quát hóa};

    % Khối 2: Group CNN (Gộp toàn bộ text vào 1 node)
    \node[draw, rounded corners, fill=orange!10, align=center, minimum width=3.5cm, minimum height=1.5cm, inner sep=0.2cm, right=1.2cm of arrow] (gcnn) {
        Bộ lọc trượt theo\\
        \textbf{mọi phần tử} của $G$\\[1ex]
        \textcolor{orange!70!black}{\footnotesize Tương biến với $G$}
    };
    \node[above=0.2cm of gcnn, font=\small\bfseries] {Tích chập nhóm};

\end{tikzpicture}
\caption{So sánh tích chập thông thường và tích chập nhóm. Tích chập nhóm là sự tổng quát hóa trực tiếp: thay vì chỉ trượt theo phép tịnh tiến, bộ lọc được trượt qua toàn bộ các phần tử của nhóm $G$.}
\label{fig:group-conv-compare}
\end{figure}

\begin{itemize}
    \item Nếu $G$ là nhóm tịnh tiến trên $\mathbb{Z}^2$, tích chập nhóm hoàn toàn trùng với CNN thông thường.
    \item Nếu $G$ bổ sung thêm các phép xoay $90°$ và phản chiếu, bộ lọc sẽ tự động bất biến với 8 phép đối xứng của hình vuông (nhóm dihedral $D_4$).
    \item Tích chập nhóm là nền tảng của \textbf{mạng tích chập tương biến nhóm}, mở rộng CNN sang các nhóm đối xứng tùy ý.
\end{itemize}

% ──────────────────────────────────────────────────────────────

\subsubsection{Lý thuyết bất biến và tham số hóa qua đại lượng bất biến}
\label{subsubsec:invariant-theory}

\textbf{Ý tưởng cốt lõi.} Thay vì thiết kế lại toàn bộ kiến trúc, ta có thể tham số hóa trực tiếp các hàm bất biến thông qua một tập hợp nhỏ các đại lượng vô hướng bất biến đơn giản.

\begin{theorem}[Định lý cơ bản cho nhóm $O(d)$]
Mọi hàm $f: (\mathbb{R}^d)^n \to \mathbb{R}$ bất biến với nhóm xoay và phản chiếu $O(d)$ đều có thể biểu diễn dưới dạng:
\begin{equation}
    f(v_1, \ldots, v_n) = h\!\left(\left\{v_i^\top v_j\right\}_{i,j=1}^{n}\right)
    \label{eq:first-fundamental-Od}
\end{equation}
trong đó $h$ là một hàm bất kỳ tác động lên ma trận Gram của các tích vô hướng.
\end{theorem}

\begin{figure}[H]
\centering
\begin{tikzpicture}[>=Stealth, thick]
    % Điểm gốc cho các vector
    \node at (0,0) [circle, fill=black, inner sep=2pt] (origin) {};

    % Vẽ các vector đầu vào
    \foreach \i/\ang/\lbl in {1/30/$v_1$, 2/100/$v_2$, 3/200/$v_3$} {
        \draw[->, blue!70!black, very thick] (origin) -- ({1.5*cos(\ang)},{1.5*sin(\ang)});
        \node[blue!70!black] at ({1.8*cos(\ang)},{1.8*sin(\ang)}) {\small \lbl};
    }

    % Dấu mũi tên chuyển đổi 1
    \node[right=1.5cm of origin] (arrow1) {\Large $\xrightarrow{\text{trích xuất}}$};

    % Khối Ma trận Gram (Gộp ma trận và chú thích vào chung 1 node để tự căn giữa)
    \node[right=0.5cm of arrow1, align=center] (gram) {
        $\begin{pmatrix} 
        v_1^\top v_1 & v_1^\top v_2 & v_1^\top v_3 \\ 
        v_2^\top v_1 & v_2^\top v_2 & v_2^\top v_3 \\ 
        v_3^\top v_1 & v_3^\top v_2 & v_3^\top v_3 
        \end{pmatrix}$\\[1.5ex]
        \textcolor{gray}{\footnotesize Ma trận Gram bất biến với $O(d)$}
    };

    % Dấu mũi tên chuyển đổi 2
    \node[right=0.5cm of gram] (arrow2) {\Large $\xrightarrow{h(\cdot)}$};

    % Đầu ra
    \node[draw, rounded corners, fill=green!20, right=0.5cm of arrow2] (output) {$f(v_1,v_2,v_3)$};

\end{tikzpicture}
\caption{Minh họa định lý cơ bản: mọi hàm bất biến với $O(d)$ đều có thể được tính từ ma trận Gram của các tích vô hướng, bất kể số chiều $d$ hay số điểm $n$.}
\label{fig:gram-matrix}
\end{figure}

\begin{itemize}
    \item \textbf{Tại sao đúng:} phép xoay và phản chiếu bảo toàn góc và khoảng cách, tức là bảo toàn chính xác các tích vô hướng $v_i^\top v_j$.
    \item \textbf{Tại sao mạnh:} định lý đảm bảo tính phổ quát, không có thông tin bất biến nào bị bỏ sót.
    \item \textbf{Với hàm tương biến $O(d)$:} mở rộng thêm tổng tuyến tính của các vector: $f(v_1,\ldots,v_n) = h(\{v_i^\top v_j\}) \cdot \sum_t v_t$.
    \item \textbf{Với tính bất biến tịnh tiến:} chỉ cần chuyển sang hệ tọa độ tương đối: $f(v_1,\ldots,v_n) = h(v_2{-}v_1, v_3{-}v_1, \ldots, v_n{-}v_1)$.
    \item \textbf{Thực tiễn:} hàm $h$ được tham số hóa bởi một mạng nơ-ron thông thường đặt trên đầu các đặc trưng bất biến đã tính sẵn.
\end{itemize}

% ──────────────────────────────────────────────────────────────

\subsubsection{Phương pháp tensor và giới hạn của bổ đề Schur}
\label{subsubsec:tensor-methods}

\textbf{Vấn đề: tại sao lớp tuyến tính tương biến đơn thuần không đủ?}

\begin{itemize}
    \item Mạng nhiều lớp tuyến tính tương biến xen kẽ hàm kích hoạt đặt ra câu hỏi: liệu có đủ sức biểu diễn để xấp xỉ mọi hàm tương biến liên tục không?
    \item \textbf{Bổ đề Schur} từ lý thuyết biểu diễn nhóm chỉ ra: với một số nhóm liên tục, các ánh xạ tuyến tính tương biến bị hạn chế nghiêm trọng.
\end{itemize}

\begin{example}
Với nhóm xoay $SO(2)$ tác động trên $\mathbb{R}^2$, ánh xạ tuyến tính tương biến duy nhất là bội của ma trận đơn vị: $L(x) = \lambda x$, $\lambda \in \mathbb{R}$. Lớp tuyến tính tương biến không thể tạo ra bất kỳ biểu diễn phong phú nào.
\end{example}

\textbf{Giải pháp: nâng đầu vào lên không gian tensor.}

Thay vì làm việc trực tiếp với $x \in \mathbb{R}^d$, ta dùng lũy thừa tensor bậc $k$:

\begin{equation}
    x^{\otimes k} = \underbrace{x \otimes x \otimes \cdots \otimes x}_{k \text{ lần}} \in (\mathbb{R}^d)^{\otimes k}
    \label{eq:tensor-power}
\end{equation}

\begin{figure}[H]
\centering
\begin{tikzpicture}[>=Stealth, thick]
    % Khai báo các khối chính với khoảng cách mép cố định (right=1cm of...)
    \node[draw, fill=blue!10, rounded corners, minimum height=0.7cm] (x) {$x \in \mathbb{R}^d$};
    
    \node[draw, fill=purple!15, rounded corners, minimum height=0.7cm, right=1cm of x] (xk) {$x^{\otimes k} \in (\mathbb{R}^d)^{\otimes k}$};
    
    \node[draw, fill=orange!20, rounded corners, minimum height=0.7cm, right=1cm of xk] (mlp) {MLP tương biến};
    
    \node[draw, fill=green!20, rounded corners, minimum height=0.7cm, right=1cm of mlp] (out) {Đầu ra tương biến};

    % Vẽ các mũi tên
    \draw[->] (x) -- node[above, font=\scriptsize] {$\otimes k$} (xk);
    \draw[->] (xk) -- (mlp);
    \draw[->] (mlp) -- (out);

    % Ghi chú bên dưới (Dùng align=center và \\ để ngắt dòng cho gọn)
    \node[below=0.2cm of xk, font=\scriptsize, text=purple!70!black, align=center] {không gian biểu diễn\\phong phú};
    
    \node[below=0.2cm of mlp, font=\scriptsize, text=orange!70!black, align=center] {universality với\\$k$ đủ lớn};
\end{tikzpicture}
\caption{Phương pháp tensor: nâng đầu vào lên không gian tensor bậc $k$ để vượt qua giới hạn của bổ đề Schur và đạt được tính phổ quát.}
\label{fig:tensor-method}
\end{figure}

\begin{itemize}
    \item Không gian tensor $(\mathbb{R}^d)^{\otimes k}$ có cấu trúc biểu diễn nhóm vô cùng phong phú, vượt qua giới hạn của bổ đề Schur.
    \item Xây dựng MLP tương biến trên $x^{\otimes k}$ đạt được tính phổ quát khi $k$ đủ lớn.
    \item \textbf{Thách thức kỹ thuật:} phân rã biểu diễn nhóm trên $(\mathbb{R}^d)^{\otimes k}$ đòi hỏi tính toán \textbf{hệ số Clebsch-Gordan}.
        \begin{itemize}
            \item Với nhóm xoay $SO(d)$ và phản chiếu $O(d)$: giải được hiệu quả.
            \item Với các nhóm hữu hạn tổng quát: là bài toán khó nổi tiếng trong lý thuyết biểu diễn.
        \end{itemize}
\end{itemize}

% ──────────────────────────────────────────────────────────────

\subsubsection{Ứng dụng: mạng đồ thị tương biến và bài toán hệ nguyên tử}
\label{subsubsec:equivariant-gnn}

\textbf{Bài toán và cấu trúc dữ liệu.}

\begin{itemize}
    \item Một hệ nguyên tử gồm $N$ nguyên tử được biểu diễn dưới dạng một đồ thị có định hướng.
    \item Mỗi nút $i$ mang: đặc trưng loại nguyên tử $h_i \in \mathbb{R}^c$ (vô hướng) và tọa độ không gian $\mathbf{r}_i \in \mathbb{R}^3$ (vector).
    \item Bài toán yêu cầu dự đoán năng lượng toàn phần $E$ của hệ, vốn là một hàm bất biến với nhóm Euclide $E(3)$ gồm phép xoay, phản chiếu và tịnh tiến, đồng thời bất biến với mọi phép hoán vị nguyên tử cùng loại.
\end{itemize}

\begin{figure}[H]
\centering
\begin{tikzpicture}[>=Stealth, thick]
    % Phân tử H2O minh họa
    \node[circle, draw, fill=red!30, minimum size=0.9cm, font=\small\bfseries] (O) at (0,0) {O};
    \node[circle, draw, fill=blue!20, minimum size=0.7cm, font=\small\bfseries] (H1) at (-1.4, -1.0) {H};
    \node[circle, draw, fill=blue!20, minimum size=0.7cm, font=\small\bfseries] (H2) at (1.4, -1.0) {H};
    \draw[-, very thick] (O) -- (H1);
    \draw[-, very thick] (O) -- (H2);

    % Nhãn đặc trưng
    \node[font=\scriptsize, text=gray, right=0.1cm of O] {$h_O,\, \mathbf{r}_O$};
    \node[font=\scriptsize, text=gray, below=0.05cm of H1] {$h_H,\, \mathbf{r}_{H_1}$};
    \node[font=\scriptsize, text=gray, below=0.05cm of H2] {$h_H,\, \mathbf{r}_{H_2}$};

    % Mũi tên ra
    \node at (3.5, -0.5) {\Large $\xrightarrow{\text{GNN tương biến}}$};

    % Đầu ra
    \node[draw, rounded corners, fill=green!20, minimum width=2.2cm, minimum height=0.9cm, font=\small] at (7, -0.5) {$E$ (năng lượng)};
    \node[font=\scriptsize, text=gray] at (7, -1.2) {bất biến với $E(3)$};
\end{tikzpicture}
\caption{Minh họa bài toán hệ nguyên tử: phân tử $\text{H}_2\text{O}$ được biểu diễn dưới dạng đồ thị. Mỗi nút mang đặc trưng vô hướng $h_i$ (loại nguyên tử) và tọa độ vector $\mathbf{r}_i$. Mô hình cần dự đoán năng lượng $E$ bất biến với mọi phép biến đổi $E(3)$.}
\label{fig:molecule-graph}
\end{figure}

\textbf{Cơ chế truyền thông tin tương biến.}

Mạng nơ-ron tương biến cập nhật lặp lại cả đặc trưng ẩn $h_i^{(\ell)}$ lẫn tọa độ $\mathbf{x}_i^{(\ell)}$ tại mỗi lớp $\ell$, đảm bảo toàn bộ quá trình tương biến với $E(3)$:

\begin{equation}
    \mathbf{x}_i^{(\ell+1)} = \mathbf{x}_i^{(\ell)}
    + C \sum_{j \neq i}
    \left(\mathbf{x}_i^{(\ell)} - \mathbf{x}_j^{(\ell)}\right)
    \cdot \varphi\!\left(h_i^{(\ell)},\; h_j^{(\ell)},\;
    \left\|\mathbf{x}_i^{(\ell)} - \mathbf{x}_j^{(\ell)}\right\|^2,\;
    a_{ij}\right)
    \label{eq:egnn-update}
\end{equation}

trong đó $\varphi$ là một mạng nơ-ron nhỏ và $a_{ij}$ là đặc trưng cạnh bổ sung.

\begin{figure}[H]
\centering
\begin{tikzpicture}[>=Stealth, thick, node distance=2cm]
    % Lớp đầu vào
    \node[circle, draw, fill=blue!20, minimum size=0.8cm, font=\scriptsize] (i_in) at (0, 1.2) {$i$};
    \node[circle, draw, fill=blue!20, minimum size=0.8cm, font=\scriptsize] (j_in) at (0, -0.8) {$j$};
    \node[font=\scriptsize, left=0.1cm of i_in] {$h_i^{(\ell)},\, \mathbf{x}_i^{(\ell)}$};
    \node[font=\scriptsize, left=0.1cm of j_in] {$h_j^{(\ell)},\, \mathbf{x}_j^{(\ell)}$};

    % Hộp message
    \node[draw, rounded corners, fill=orange!15, minimum width=2.5cm, minimum height=1.2cm,
          font=\scriptsize, align=center] (msg) at (3.5, 0.2)
          {Tính message \\ $\varphi(h_i, h_j, \|\mathbf{x}_i{-}\mathbf{x}_j\|^2, a_{ij})$};

    % Lớp đầu ra
    \node[circle, draw, fill=green!30, minimum size=0.8cm, font=\scriptsize] (i_out) at (7, 0.2) {$i$};
    \node[font=\scriptsize, right=0.1cm of i_out] {$h_i^{(\ell+1)},\, \mathbf{x}_i^{(\ell+1)}$};

    % Mũi tên
    \draw[->] (i_in) -- (msg);
    \draw[->] (j_in) -- (msg);
    \draw[->] (msg) -- (i_out) node[midway, above, font=\scriptsize] {cập nhật};

    % Nhãn đặc tính
    \node[font=\scriptsize, text=orange!70!black, below=0.1cm of msg]
        {chỉ dùng $\|{\cdot}\|^2$ $\Rightarrow$ bất biến với xoay};
\end{tikzpicture}
\caption{Cơ chế cập nhật tương biến: hàm message $\varphi$ chỉ nhận khoảng cách bình phương $\|\mathbf{x}_i - \mathbf{x}_j\|^2$ (một đại lượng bất biến) thay vì tọa độ tuyệt đối, đảm bảo kết quả tương biến với nhóm $E(3)$.}
\label{fig:egnn-message}
\end{figure}

\begin{itemize}
    \item Hàm $\varphi$ chỉ nhận \textbf{khoảng cách bình phương} $\|\mathbf{x}_i - \mathbf{x}_j\|^2$, một đại lượng bất biến với $E(3)$, thay vì tọa độ tuyệt đối.
    \item Cập nhật tọa độ theo hướng $(\mathbf{x}_i - \mathbf{x}_j)$ đảm bảo đầu ra tương biến: xoay đầu vào thì đầu ra cũng xoay theo đúng cách đó.
    \item Đặc trưng vô hướng $h_i$ và tọa độ vector $\mathbf{x}_i$ được cập nhật đồng thời qua nhiều lớp, làm giàu dần biểu diễn.
\end{itemize}

\textbf{Lợi ích đặc biệt: tính toán lực miễn phí qua vi phân tự động.}

Sau khi mô hình dự đoán được năng lượng $E$, lực tác dụng lên từng nguyên tử thu được hoàn toàn miễn phí:

\begin{equation}
    \mathbf{F}_i = -\frac{\partial E}{\partial \mathbf{r}_i}
    \label{eq:force-autograd}
\end{equation}

\begin{figure}[H]
\centering
\begin{tikzpicture}[>=Stealth, thick]
    % Khối GNN gốc
    \node[draw, rounded corners, fill=blue!10, minimum height=0.8cm, font=\small] 
        (model) {GNN tương biến};

    % Khối Năng lượng (đặt thẳng sang bên phải khối GNN)
    \node[draw, rounded corners, fill=green!20, minimum height=0.7cm, font=\small, right=2.5cm of model] 
        (E) {$E$ (năng lượng)};

    % Khối Lực (đặt thẳng ngay bên dưới khối Năng lượng)
    \node[draw, rounded corners, fill=orange!20, minimum height=0.7cm, font=\small, below=1cm of E] 
        (F) {$\mathbf{F}_i = -\frac{\partial E}{\partial \mathbf{r}_i}$};

    % Các đường mũi tên liên kết
    \draw[->] (model) -- node[midway, above, font=\scriptsize] {dự đoán trực tiếp} (E);
    \draw[->, dashed] (E) -- node[midway, right, font=\scriptsize] {autograd} (F);
    \draw[->, dashed] (model) -- node[midway, below, sloped, font=\scriptsize, text=gray] {(miễn phí)} (F);
\end{tikzpicture}
\caption{Sơ đồ tính lực qua vi phân tự động: mô hình chỉ được huấn luyện dự đoán năng lượng $E$ (vô hướng), nhưng đồng thời cung cấp $3N$ thành phần lực vector nhất quán và đảm bảo tính đối xứng.}
\label{fig:force-autograd}
\end{figure}

\begin{itemize}
    \item Mô hình chỉ cần học dự đoán một đại lượng vô hướng (năng lượng), nhưng đồng thời cung cấp dự đoán cho $3N$ thành phần lực vector.
    \item Các lực này \textbf{nhất quán} với nhau và với năng lượng theo định luật bảo toàn, vì chúng đến từ cùng một hàm thế năng.
    \item Đây là lợi thế quyết định so với các mô hình dự đoán lực trực tiếp, vốn không đảm bảo tính nhất quán.
\end{itemize}

\textbf{Ứng dụng mở rộng.} Nguyên lý tương biến $E(3)$ trong mạng đồ thị không chỉ giới hạn ở hệ nguyên tử, mà còn được áp dụng rộng rãi trong:

\begin{itemize}
    \item \textbf{Sinh protein:} các mô hình Chroma và RFDiffusion dùng GNN tương biến kết hợp với hệ tọa độ địa phương để tạo sinh cấu trúc protein.
    \item \textbf{Dự đoán liên kết phân tử:} tương biến $SE(3)$ cho phép mô hình nắm bắt hướng và hình dạng của túi liên kết kháng nguyên-kháng thể.
    \item \textbf{Điều khiển robot:} tương biến $SO(2)$/$SE(3)$ trong học tăng cường cho phép chính sách tối ưu tự động điều chỉnh khi robot bị xoay hoặc dịch chuyển trong không gian (được trình bày chi tiết ở Mục~\ref{subsec:robotics}).
\end{itemize}

% ─────────────────────────────────────────────────────────────
% ỨNG DỤNG & GIỚI HẠN
% ─────────────────────────────────────────────────────────────

\subsection{Ứng dụng: Equivariance trong Robotics}
\label{subsec:robotics}
% Trình bày:
%   - Trực quan: xoay robot + action → action cũng xoay
%   - Theorem: MDPs có equivariant optimal policies:
%     π*(g·state) = g·π*(state) ∀g∈G
%   - SO(2)/SE(3) symmetries trong grasp learning, manipulation
%   - Kết quả: sample-efficient learning
%   (Wang et al CoRL 2021, ICLR 2022)

\subsection{Giới hạn: Vấn đề Symmetry Breaking}
\label{subsec:symmetry-breaking}
% Trình bày:
%   - Theorem: equivariant architectures không thể break symmetry
%     Proof: G_x ⊆ G_y (output always more symmetric than input)
%   - Hệ quả: không thể generate non-symmetric output
%     từ symmetric input (problem cho generative models)
%   - Giải pháp: weight perturbation, gradients, sets
%   (Wang et al ICML 2024)

\subsection{So sánh hai hướng tiếp cận}
\label{subsec:comparison}
% Bảng tổng hợp:
%
%  | Tiêu chí         | Off-the-shelf       | Constructed        |
%  |------------------|---------------------|--------------------|
%  | Flexibility      | Cao (dùng bất kỳ)   | Thấp (thiết kế lại)|
%  | Symmetry đảm bảo | Approximate (Aug.)  | Exact              |
%  |                  | Exact (Canon/Avg.)  |                    |
%  | Expressiveness   | Kế thừa base model  | Phụ thuộc design   |
%  | Computational    | Có thể đắt (Avg.)   | Efficient          |
%  | Continuity       | Vấn đề (Canon.)     | Không vấn đề       |
%
% Khi nào dùng gì:
%   - Off-the-shelf: khi có base model tốt, symmetries ít đặc biệt
%   - Constructed: khi symmetry rõ ràng, data đắt (science/robotics)

\clearpage
